% Supplementary information, language and reviewer-audit revision 6 October 2026.
\documentclass[pdflatex,sn-vancouver-num]{sn-jnl}
\usepackage{booktabs}
\usepackage{tabularx}
\usepackage{array}
\usepackage{amsmath,amssymb}
\usepackage{url}
\usepackage{graphicx}
\usepackage{seqsplit}
\usepackage{longtable}
\usepackage{placeins}
\raggedbottom
\unnumbered

\begin{document}
\renewcommand{\thetable}{S\arabic{table}}

\title{Supplementary information for: Early versus recent personal baselines in repeated psychological self-reports: a cross-dataset prediction study}
\author[1]{\fnm{Runqing} \sur{Liao}}\email{202411055222@stu.hstc.edu.cn}
\author*[1]{\fnm{Zhining} \sur{Wang}}\email{mathwangzhining@163.com}
\affil*[1]{\orgdiv{School of Mathematics and Statistics}, \orgname{Hanshan Normal University}, \orgaddress{\city{Chaozhou}, \postcode{521041}, \country{China}}}
\maketitle

This supplement is organized around the evidential sequence of the study. Supplementary Methods 1--7 document the three development datasets, the equal-information comparison, and the analyses that challenged alternative explanations. Supplementary Methods 8 documents the independently sourced Corona Health extension and its protocol lock. Tables S1--S2 describe sample flow and construct mapping; Tables S3--S20 report development-dataset diagnostics; and Tables S21--S24 report the external estimates, AR(1) model checks, and the distinction between conclusions supported by the diagnostics and mechanisms that remain unresolved.

\section{Supplementary Methods 1: cohorts, chronology, and inclusion}

All analyses used de-identified secondary data. Within each dataset--outcome pair, complete reports were ordered chronologically. The first eight complete reports formed the calibration block. A participant entered an equal-information prediction analysis only if at least one later complete target existed; participant-specific age-slope analyses additionally required at least five targets and variation in baseline age. No statistical outlier rule was used to remove observed outcome values. Scale-range validation occurred before analysis. Repeated reports were never treated as independent participants.

The study used a staged evidential chronology. CES and Marian supported initial mechanism development. The first association-blind Dejonckheere analysis compared a training-weighted early eight-report mean/last-value baseline with a cumulative online mean/last-value rule; a weighted recent-eight rule was added after those results were available. The equal-weight B8--L8 comparison reported here was defined subsequently, and all null, comparator, temporal-generalisability, machine-learning, semi-synthetic, window, observation-process, and pseudo-origin analyses in those three datasets are explicitly post hoc. On 5 October 2026, Corona Health structural fields, scale ranges, completeness, duplicate timestamps, and report counts were audited without calculating B8--L8 associations. A dated protocol then froze the primary PHQ-9 outcome, secondary GAD-7 outcome, eligibility, estimands, duplicate handling, long-gap sensitivity, ordinal checks, and stationary challenge before outcome comparison. This is a protocol-locked external extension, not preregistration. The embargoed OSF registration (\url{https://osf.io/n2j3m/}) concerned a separate CES--PSYCHE-D activity and sleep question and is not claimed as preregistration for the present study.

\section{Supplementary Methods 2: equal-information estimands and null models}

For participant $i$ and complete-report index $t>8$, the frozen and recent baselines were
\[
B8_i=\frac{1}{8}\sum_{j=1}^{8}y_{ij},\qquad
L8_{it}=\frac{1}{8}\sum_{j=t-8}^{t-1}y_{ij}.
\]
Update benefit was $U_{it}=|y_{it}-B8_i|-|y_{it}-L8_{it}|$. Baseline age was elapsed days from the eighth report in CES and scheduled-opportunity distance in the two experience-sampling datasets. Age was transformed as $\log(1+\mathrm{age})$ and divided by its dataset--outcome target-level standard deviation. For participants with sufficient targets, an ordinary least-squares slope of $U_{it}$ on within-person-centred scaled age was fitted. The arithmetic mean of participant slopes was the inferential estimand.

Time permutation randomly reordered non-missing scores within participant while retaining assessment positions and target counts. The circular-shift challenge selected a non-zero cyclic shift within participant, preserving the score multiset, missing positions, and cyclic adjacent-pair structure while relocating the study start. B8, L8, target errors, and slopes were reconstructed for every one of 500 draws. One-sided Monte Carlo values were $(1+k)/501$. Four primary outcomes used a Bonferroni threshold of 0.0125.

The principal post hoc null was a participant-matched stationary Gaussian AR(1) parametric bootstrap. For each participant and outcome, the mean, residual innovation scale, and lag-1 coefficient were estimated from the complete observed sequence. Participant coefficients were shrunk towards the outcome-level estimate and clipped to $[-0.95,0.95]$. Each simulated development sequence began from the implied stationary Gaussian distribution, advanced in complete-report index, and was clipped to the admissible scale bounds. Clipping did not preserve the observed integer score grid or exact marginal moments. Five thousand sequences were generated. Observed report counts and age covariates were retained, but transitions did not depend on elapsed time and latent scores were not generated over missing opportunities before masking. B8, L8, participant-balanced B8-minus-L8 MAE, and mean age gradients were reconstructed. One-sided Monte Carlo values were $(1+k)/5001$. The model tests the fitted Gaussian first-order process, not all stationary psychological processes. A separately frozen post hoc diagnostic used 200 additional sequences per participant to compare observed and simulated marginal features (Table S23).

An additional post hoc circular geometric stationary bootstrap sampled each participant's empirical sequence with mean block lengths of 8, 16, and 32 reports (1,000 draws each). Initial source positions were uniform; at each subsequent position the source index restarted uniformly with probability $1/L$ and otherwise advanced circularly. Conditional on the observed sequence, this preserves the bounded empirical score grid and its marginal distribution in expectation while retaining different amounts of local dependence. Block lengths were diagnostic choices rather than calibrated psychological timescales, and the procedure has no universal finite-sample calibration guarantee.

\section{Supplementary Methods 3: temporal generalisability and machine learning}

For temporal generalisability, each participant's post-calibration targets were divided chronologically into non-overlapping early and late blocks. Each block required at least five targets and age variation. Slopes and mean update benefits were estimated independently in the two blocks. Pearson and Spearman correlations were calculated across participants; Pearson intervals used 2,000 participant bootstrap samples. Standardised analyses excluded zero-variance participants and recomputed raw mean-benefit and age-gradient correlations on that same eligibility set. Both blockwise mean benefits share B8; their correlation therefore describes repeatability of a shared-reference task statistic rather than an independent stable trait.

The exploratory machine-learning target was late-block mean update benefit. Predictors were restricted to the first eight reports: mean, sample standard deviation, mean absolute successive difference, and linear trend. Ridge regression used five participant-disjoint outer folds and five-fold inner selection of $\alpha\in\{0.01,0.1,1,10,100\}$. Scaling and tuning were fitted only within training data. Each of 500 label permutations reran the complete nested procedure. The frozen gate required more than 10\% RMSE improvement, a participant-bootstrap interval below zero for relative RMSE change, positive out-of-fold correlation, and Bonferroni-adjusted one-sided $p\leq0.0125$.

\section{Supplementary Methods 4: sustained-change stress test}

The stress test retained each participant's observed residual sequence and added a directional offset to future targets. Higher values indicated worsening for the four primary outcomes. In the primary scenario, the offset began at the midpoint, increased linearly from one eighth to one pooled within-person outcome standard deviation over eight complete reports, and remained at one standard deviation for eight additional reports. Eligibility required at least 32 post-calibration targets.

For an injected offset sequence $d_{it}$, B8 did not adapt. L8's injected adaptation at target $t$ was $8^{-1}\sum_{j=t-8}^{t-1}d_{ij}$, so neither $d_{it}$ nor future offsets entered the current baseline. A50 was $0.5B8+0.5L8$. Sustained-phase signal retention was the injection-induced change in worsening-direction deviation divided by the imposed plateau shift. Because the injected trajectory and baseline weights were fixed, retention is a structural property of the schedule, not a clinical effect estimate.

For the secondary alarm comparison, each outcome--method threshold was the weighted 95th percentile of pre-injection original deviation scores, with each participant receiving equal total weight. Matched post-onset original scores provided the false-alarm rate. Injected scores provided detection rate; their difference was the excess alarm rate. The full grid crossed 0.5- and 1.0-SD amplitudes, onset fractions 0.40, 0.50, and 0.60, and ramp lengths 4, 8, and 16. Ineligible dataset--scenario combinations remained recorded in the audit rather than being silently dropped.

\section{Supplementary Methods 5: uncertainty, integrity, and reproducibility}

Participant-balanced MAE was the arithmetic mean of participant-specific MAEs. Participant-balanced RMSE was the square root of the arithmetic mean of participant-specific mean squared errors. Unless otherwise stated, intervals used 2,000 bootstrap resamples of participants. These intervals quantify sampling variation among observed participants, conditional on the selected datasets and analysis decisions. They are not intervals over hypothetical datasets. All random procedures used seed 20260918. Automated tests verified strict temporal ordering, current/future exclusion, equal eight-value construction, disjoint participant folds, nested training-only preprocessing, null-model invariants, injection ordering, participant-level resampling, and output reconciliation.

\section{Supplementary Methods 6: post hoc comparator and bounded-score analyses}

Additional predictors were computed from strictly prior reports. ``Last'' was the immediately preceding complete value. ``Cumulative'' was the mean of all complete values before the target. EWM8 was initialized at B8 and then updated as $m_t=(1-2/9)m_{t-1}+(2/9)y_{t-1}$. Median8 summarized the preceding eight reports. For bounded integer outcomes, Mode8 selected the tied modal value that appeared most recently in the eight-report history. Online AR(1) re-estimated an intercept and lag coefficient from preceding reports at each target and clipped the coefficient to the stationary range. OracleMean used the complete participant series and is a retrospective diagnostic comparator. For PHQ-2 and Marian outcomes, B8 and L8 predictions were rounded and clipped to the admissible score set before computing exact accuracy, within-one accuracy, and absolute ordinal error. Participant-SD-standardised MAE divided each participant's MAE by that participant's observed within-person standard deviation before averaging. These additions are interpreted descriptively.

\section{Supplementary Methods 7: origin, window, and observation sensitivities}

These analyses were written into dated post hoc protocols before execution. Window-length analyses compared early and recent equal-weight means for $K=4,6,8,$ and 12 on common targets beginning at complete-report index 13. CES recall separation used elapsed days between the target and immediately preceding retained PHQ-2 assessment. Each Marian target retained its scheduled-opportunity counter. Response probabilities were joined by participant, outcome, and that counter and came from five folds with no participant overlap between training and test sets. An L2 logistic model used strictly prior cumulative mean and last score, prior complete-report count, prior completion rate, current missing streak, day, beep, and counter. Regularisation and standardisation were selected within training participants; estimated response probabilities were clipped to 0.10--0.99, reciprocal inverse-probability weights were calculated, and those weights were normalised within participant. EMA timing analyses either removed all day-1 observations and re-established the first eight complete reports or subtracted prompt-position means estimated from all complete reports. The latter is a retrospective descriptive adjustment, not a strictly past-only prediction procedure. Pseudo-origin placement used the retrospectively observed complete-series length; consequently, repeated-origin results diagnose dependence on the original study start but do not represent prospective rolling-origin validation.

For repeated-origin analyses, a participant required at least 21 complete observations. Early, middle, and late origins were the first eligible index, the midpoint of the eligible starting-index range, and the final index retaining at least five evaluation targets. Eight reports formed the fixed baseline, the next eight reports were a washout, and evaluation began at the seventeenth report after each origin. B8 and L8 therefore shared no calibration values at the first evaluated target. Origins were not independent because they reused participants and sometimes observations. All contrasts and slope intervals used 2,000 participant bootstrap resamples with seed 20260918.

\section{Supplementary Methods 8: protocol-locked Corona Health extension}

Corona Health version 2 was obtained from B2SHARE (doi:10.23728/b2share.cgf63-kme28). The EMA file SHA-256 was \texttt{\seqsplit{FBFF34486AC2D4FF61D5E92B2CAA7276CEC51719973FE232DF3096D3357F8D87}}; the codebook SHA-256 was \texttt{\seqsplit{551B65FDB308EDBA9809B2B78FD30D99EED8AC9B6964BBC8594400E256CA96B1}}. The CSV contained unquoted comma-separated values in fields after the retained scale items. Parsing therefore addressed columns by their header positions and stopped after the GAD-7 items; every retained row was required to contain all addressed positions. Later free-text, device, and multiselect fields were not used.

PHQ-9 was the frozen primary outcome and GAD-7 the frozen secondary outcome. Totals required every item, each within 0--3. The app instructed participants to report symptoms over the previous week. Records were ordered by timestamp. Forty-seven repeated participant--timestamp groups had identical PHQ-9 item values and were collapsed; no conflicting complete PHQ-9 duplicates were found. Prediction eligibility required nine complete reports. Slope eligibility required at least five targets and age variation. Baseline age was elapsed days from the eighth report.

The primary participant-balanced MAE contrast, age gradient, and 2,000-resample participant bootstrap followed Supplementary Methods 2 and 5. In the long-gap sensitivity, a gap longer than 30 days started a new segment; B8 and L8 were reconstructed within each segment, and only segments with at least nine records contributed. Rounded predictions were clipped to 0--27 for PHQ-9 and 0--21 for GAD-7. Exact accuracy, within-two-point accuracy, and ordinal MAE were first averaged within participant and then across participants. Median8 was a frozen descriptive comparator.

The PHQ-9 stationary challenge generated 5,000 participant-matched Gaussian AR(1) sequences in complete-report index. Participant lag-1 estimates were shrunk towards the pooled estimate and clipped to $[-0.95,0.95]$; observed sequence lengths, means, and SDs parameterised the generator but were not forced to match exactly in each simulated sequence. The Corona generator was not clipped to the PHQ-9 scale bounds, a choice exposed by the marginal model check in Table S23. Because the source schedule was irregular, this null does not represent a continuous-time process and does not reproduce elapsed gaps. The protocol, script, aggregate outputs, source data for Fig.~2 of the main manuscript, and run manifest are included in the analysis repository.

\begin{table*}[!htbp]
\caption{Cohort flow and analytic roles.}\label{tab:s1}
\centering\scriptsize
\setlength{\tabcolsep}{3pt}
\begin{tabularx}{\textwidth}{@{}l>{\raggedright\arraybackslash}Xrrrrrr@{}}
\toprule
Dataset & Outcome role and schedule & \shortstack{Source\\IDs} & \shortstack{Source\\rows} & \shortstack{Complete\\records} & \shortstack{Eligible\\IDs} & \shortstack{Post-calibration\\targets} & \shortstack{Slope\\IDs}\\
\midrule
Dejonckheere & Five 0--100 momentary items; 14 days & 100 & 9,800 & 8,702 & 100 & 7,902 per outcome & 100\\
CES & PHQ-2 (0--6); up to four years & 110 & 5,894 & 5,407 & 105 & 4,541 & 99\\
Marian & Two 1--4 four-hour items; 21 days & 145 & 9,135 & 8,231 & 145 & 7,071 per outcome & 145\\
Corona Health & PHQ-9 (0--27); irregular, one-week recall & 1,487 & 11,541 & 11,432 & 324 & 5,618 & 232\\
Corona Health & GAD-7 (0--21); irregular, one-week recall & 1,488 & 11,541 & 11,494 & 327 & 5,663 & 232\\
\bottomrule
\end{tabularx}
\begin{flushleft}\footnotesize Eligibility required at least nine retained reports (eight calibration reports plus one target). Slope analyses additionally required at least five targets and age variation. Corona Health source participants count identifiers with a complete retained scale; complete records are after collapse of identical participant--timestamp duplicates. The supplied CES development-derived table contained 5,894 visits; panel construction retained 5,407. This is not a scheduled-opportunity denominator, and preprocessing upstream of the supplied derived table was unavailable for audit. Stress-test eligibility was scenario-specific and is reported in Table S6.\end{flushleft}
\end{table*}

\begin{table*}[!htbp]
\caption{Variable and construct mapping.}\label{tab:s2}
\centering\footnotesize
\begin{tabularx}{\textwidth}{@{}ll>{\raggedright\arraybackslash}X>{\raggedright\arraybackslash}X@{}}
\toprule
Dataset & Variable & Scale / reference period & Interpretation boundary\\
\midrule
Dejonckheere & sad, stressed, happy, relaxed, angry & 0--100; momentary & Single affect/stress items; not diagnoses\\
CES & PHQ-2 & 0--6; previous two weeks & Depression symptom screener; not a diagnostic outcome\\
Marian & depressed, anhedonia & 1--4; preceding four hours & Depression-related states; not interchangeable with standard PHQ-2 or ``right now'' ratings\\
Corona Health & PHQ-9 item sum & 0--27; previous week & Adapted recall period; symptom score, not diagnosis\\
Corona Health & GAD-7 item sum & 0--21; previous week & Secondary construct-generalisation outcome; not diagnosis\\
All & B8 & Mean of first 8 complete reports & Frozen personal reference\\
All & L8 & Mean of 8 most recent strictly prior complete reports & Time-local personal reference\\
All & A50 & $0.5B8+0.5L8$ & Post-hoc anchored stress-test comparator\\
\bottomrule
\end{tabularx}
\end{table*}

\begin{table*}[!htbp]
\caption{Circular-shift robustness across all outcomes.}\label{tab:s3}
\centering\small
\resizebox{\textwidth}{!}{%
\begin{tabular}{@{}llrrrr@{}}
\toprule
Dataset & Outcome & Observed slope & Null mean & Null 95\% range & One-sided $p$\\
\midrule
Dejonckheere & Sadness & 0.533 & 0.095 & $-0.216$, 0.434 & 0.005988\\
Dejonckheere & Stress & 1.568 & 0.166 & $-0.223$, 0.542 & 0.001996\\
Dejonckheere & Happiness & 1.275 & 0.224 & $-0.050$, 0.538 & 0.001996\\
Dejonckheere & Relaxation & 1.194 & 0.241 & $-0.091$, 0.573 & 0.001996\\
Dejonckheere & Anger & 0.523 & 0.062 & $-0.169$, 0.335 & 0.001996\\
CES & PHQ-2 & 0.110 & 0.019 & $-0.010$, 0.051 & 0.001996\\
Marian & Depressed mood & 0.059 & 0.005 & $-0.008$, 0.017 & 0.001996\\
Marian & Anhedonia & 0.052 & 0.008 & $-0.007$, 0.023 & 0.001996\\
\bottomrule
\end{tabular}%
}
\begin{flushleft}\footnotesize Primary family: sadness, stress, PHQ-2, and depressed mood; Bonferroni threshold 0.0125. Secondary-outcome values are descriptive support.\end{flushleft}
\end{table*}

\begin{table*}[!htbp]
\caption{Early--late temporal generalisability in primary outcomes.}\label{tab:s4}
\centering\small
\resizebox{\textwidth}{!}{%
\begin{tabular}{@{}llrrcc@{}}
\toprule
Dataset & Outcome & Participants & Quantity & Pearson $r$ (95\% CI) & Spearman $r$\\
\midrule
Dejonckheere & Sadness & 100 & Age slope & $-0.100$ ($-0.378$, 0.192) & 0.036\\
 & & 100 & Mean benefit & 0.650 (0.530, 0.763) & 0.663\\
Dejonckheere & Stress & 100 & Age slope & $-0.084$ ($-0.305$, 0.159) & 0.060\\
 & & 100 & Mean benefit & 0.496 (0.352, 0.640) & 0.568\\
CES & PHQ-2 & 97 & Age slope & $-0.279$ ($-0.447$, $-0.098$) & $-0.290$\\
 & & 97 & Mean benefit & 0.749 (0.367, 0.884) & 0.407\\
Marian & Depressed mood & 145 & Age slope & 0.086 ($-0.185$, 0.337) & 0.040\\
 & & 145 & Mean benefit & 0.690 (0.554, 0.800) & 0.663\\
\bottomrule
\end{tabular}%
}
\end{table*}

\begin{table*}[!htbp]
\caption{Exploratory nested participant-held-out ridge results.}\label{tab:s5}
\centering\small
\resizebox{\textwidth}{!}{%
\begin{tabular}{@{}llrrrrc@{}}
\toprule
Dataset & Outcome & Participants & RMSE change, \% (95\% CI) & OOF $r$ & Permutation $p$ & Gate\\
\midrule
Dejonckheere & Sadness & 100 & $-10.40$ ($-17.70$, $-0.85$) & 0.412 & 0.001996 & Pass\\
Dejonckheere & Stress & 100 & $-0.41$ ($-4.54$, 5.01) & 0.138 & 0.153693 & Fail\\
CES & PHQ-2 & 97 & $-0.50$ ($-2.87$, 2.61) & 0.066 & 0.113772 & Fail\\
Marian & Depressed mood & 145 & $-14.48$ ($-22.19$, $-3.71$) & 0.516 & 0.001996 & Pass\\
\bottomrule
\end{tabular}%
}
\begin{flushleft}\footnotesize OOF, out-of-fold. Negative RMSE change favours ridge over the training-fold mean. Every primary permutation reran inner tuning and outer fitting.\end{flushleft}
\end{table*}

\begin{table*}[!htbp]
\caption{Primary sustained-change stress-test results.}\label{tab:s6}
\centering\small
\resizebox{\textwidth}{!}{%
\begin{tabular}{@{}llcrrrrr@{}}
\toprule
Dataset & Outcome & Method & Participants & Signal retained & Original MAE gain, \% & Excess alarm, points (95\% CI)\\
\midrule
Dejonckheere & Sadness & B8 & 100 & 1.000 & 0.00 & 3.19 (2.19, 4.25)\\
 & & A50 & 100 & 0.582 & 8.79 & 2.88 (2.00, 3.88)\\
 & & L8 & 100 & 0.164 & 9.97 & 1.19 (0.69, 1.75)\\
Dejonckheere & Stress & B8 & 100 & 1.000 & 0.00 & 10.06 (7.44, 13.00)\\
 & & A50 & 100 & 0.582 & 12.78 & 3.94 (2.69, 5.13)\\
 & & L8 & 100 & 0.164 & 18.27 & 2.06 (1.31, 2.88)\\
CES & PHQ-2 & B8 & 71 & 1.000 & 0.00 & 5.28 (2.99, 7.83)\\
 & & A50 & 71 & 0.582 & 15.73 & 5.02 (2.46, 8.28)\\
 & & L8 & 71 & 0.164 & 23.38 & 1.94 (1.06, 2.90)\\
Marian & Depressed mood & B8 & 141 & 1.000 & 0.00 & 5.14 (3.50, 7.05)\\
 & & A50 & 141 & 0.582 & 16.42 & 3.41 (2.48, 4.43)\\
 & & L8 & 141 & 0.164 & 27.95 & 1.91 (1.33, 2.53)\\
\bottomrule
\end{tabular}%
}
\begin{flushleft}\footnotesize Primary scenario: one pooled within-person outcome SD, midpoint onset, eight-report ramp, and eight-report plateau. Scenario participant counts refer to records eligible for the specified injection schedule, whereas original MAE gains were calculated in the full prediction sample; these columns therefore have different denominators. Signal-retention values are deterministic under the imposed trajectory and fixed baseline weights. Excess alarm is injected detection rate minus matched-original false-alarm rate; intervals use participant bootstrap resampling. Thresholds are methodological, not clinical, and no observed deterioration labels were available.\end{flushleft}
\end{table*}

\begin{table*}[!htbp]
\caption{Signal-retention sensitivity by ramp length.}\label{tab:s7}
\centering\small
\begin{tabular}{@{}lccc@{}}
\toprule
Ramp length & B8 & A50 & L8\\
\midrule
4 reports & 1.000 & 0.645 & 0.289\\
8 reports & 1.000 & 0.582 & 0.164\\
16 reports & 1.000 & 0.541 & 0.082\\
\bottomrule
\end{tabular}
\begin{flushleft}\footnotesize Values were identical across amplitudes, eligible onset positions, and outcomes because retention follows from the fixed update rule and ramp geometry. Of 144 prespecified dataset--outcome--scenario combinations, 140 were analysed and four were recorded as ineligible because no participant retained the required post-onset horizon. Complete data-dependent alarm summaries are supplied as source data.\end{flushleft}
\end{table*}

\begin{table*}[!htbp]
\caption{Reproducibility record.}\label{tab:s8}
\centering\small
\begin{tabularx}{\textwidth}{@{}lX@{}}
\toprule
Item & Recorded value\\
\midrule
Environment & Python 3.12.3; NumPy 2.3.5; pandas 2.3.3; scikit-learn 1.7.2; SciPy 1.17.1 where used\\
Random seed & 20260918\\
Equal-weight input SHA-256 & \texttt{\seqsplit{21d0a8c0abf78492d7f87bc2e26208403e07b88a1ec27bac868148366a8717d9}}\\
Equal-weight protocol SHA-256 & \texttt{\seqsplit{d1249668b8aeb20af76a971cb1425f853d5043fa28a609720d1d2e10d1b448db}}\\
Circular-shift protocol SHA-256 & \texttt{\seqsplit{7f13ca7ca451cb9abbee8c732dbb53cfbda02f4847bbd2a38b18bfd2148a2fca}}\\
Temporal-generalisability protocol SHA-256 & \texttt{\seqsplit{ed386a1f42d4e61502c6d9a4e77a00dc8a8f35e1e945298c265eab04da89ec9b}}\\
Machine-learning protocol SHA-256 & \texttt{\seqsplit{c0756f4d019f07b441de7b4f4d7030f92ce355986796cdef9e296c4e87b99403}}\\
Stress-test protocol SHA-256 & \texttt{\seqsplit{2c682db3a1cdc5570d23a4b5691bc9bab6755ff6c1d55834d50d731aab7ae1ae}}\\
Stress-test script SHA-256 & \texttt{\seqsplit{7b145585f4056a47284054e1655beac2f25d968b8d7b341f131fde619905281a}}\\
\bottomrule
\end{tabularx}
\begin{flushleft}\footnotesize Hashes identify the frozen analysis artifacts used for this draft. The synchronized code and protocol package is available as GitHub release v1.3.1 at \url{https://github.com/liaorunqing/psychology/releases/tag/v1.3.1} and is permanently archived in Zenodo at \url{https://doi.org/10.5281/zenodo.23216022}.\end{flushleft}
\end{table*}

\begin{table*}[!htbp]
\caption{Absolute participant-balanced error for leakage-safe baseline rules.}\label{tab:s9}
\centering\small
\resizebox{\textwidth}{!}{%
\begin{tabular}{@{}llccccc@{}}
\toprule
Dataset & Outcome & B8 & Last & Cumulative & L8 & EWM8\\
\midrule
Dejonckheere & Sadness & 11.786/18.489 & 11.115/20.795 & 10.797/16.850 & 10.612/17.152 & 10.282/16.711\\
Dejonckheere & Stress & 21.094/26.673 & 18.525/27.929 & 18.596/23.358 & 17.239/23.167 & 16.853/22.652\\
Dejonckheere & Happiness & 17.821/23.034 & 16.411/23.383 & 15.125/19.835 & 14.489/19.522 & 14.106/19.055\\
Dejonckheere & Relaxation & 20.996/26.312 & 19.577/27.584 & 18.126/23.007 & 17.254/22.750 & 16.859/22.280\\
Dejonckheere & Anger & 10.670/17.323 & 10.559/20.355 & 9.730/15.771 & 9.580/16.103 & 9.388/15.886\\
CES & PHQ-2 & 0.975/1.437 & 0.722/1.248 & 0.814/1.178 & 0.747/1.120 & 0.724/1.079\\
Marian & Depressed mood & 0.526/0.754 & 0.335/0.697 & 0.437/0.635 & 0.379/0.619 & 0.367/0.591\\
Marian & Anhedonia & 0.567/0.795 & 0.360/0.741 & 0.460/0.667 & 0.398/0.641 & 0.388/0.621\\
\bottomrule
\end{tabular}%
}
\begin{flushleft}\footnotesize Each cell is MAE/RMSE. MAE is the arithmetic mean of participant-specific MAEs; RMSE is the square root of the arithmetic mean of participant-specific mean squared errors. Main Table 1 provides B8 and L8 MAE intervals. Main Table 2 reports comparator point estimates; participant-bootstrap intervals for comparator contrasts are available in the archived source table and are not reproduced in that main-text table. No cross-scale pooling was performed.\end{flushleft}
\end{table*}

\begin{table*}[!htbp]
\caption{Bounded-score sensitivity: L8 minus B8.}\label{tab:s10}
\centering\small
\resizebox{\textwidth}{!}{%
\begin{tabular}{@{}llrccc@{}}
\toprule
Dataset & Outcome & Participants & Exact accuracy difference (95\% CI) & Within-one difference (95\% CI) & Ordinal MAE difference (95\% CI)\\
\midrule
CES & PHQ-2 & 105 & 0.095 (0.050, 0.141) & 0.076 (0.041, 0.113) & $-0.248$ ($-0.350$, $-0.162$)\\
Marian & Depressed mood & 145 & 0.157 (0.110, 0.204) & 0.020 (0.005, 0.038) & $-0.177$ ($-0.228$, $-0.127$)\\
Marian & Anhedonia & 145 & 0.148 (0.103, 0.195) & 0.022 (0.005, 0.044) & $-0.171$ ($-0.228$, $-0.120$)\\
\bottomrule
\end{tabular}%
}
\begin{flushleft}\footnotesize Predictions were rounded and clipped to each outcome's admissible score set. Positive accuracy differences and negative MAE differences favour L8. Intervals use 2,000 participant bootstrap resamples.\end{flushleft}
\end{table*}

\begin{table*}[!htbp]
\caption{Analytic target distributions and scale-boundary frequencies.}\label{tab:s11}
\centering\small
\resizebox{\textwidth}{!}{%
\begin{tabular}{@{}llrrccrr@{}}
\toprule
Dataset & Outcome & Participants & Targets & Mean (SD) & Range & Floor, \% & Ceiling, \%\\
\midrule
Dejonckheere & Sadness & 100 & 7,902 & 11.57 (18.65) & 0--100 & 44.70 & 0.28\\
Dejonckheere & Stress & 100 & 7,902 & 25.93 (26.48) & 0--100 & 24.92 & 0.72\\
Dejonckheere & Happiness & 100 & 7,902 & 62.88 (22.69) & 0--100 & 1.04 & 8.01\\
Dejonckheere & Relaxation & 100 & 7,902 & 62.17 (24.98) & 0--100 & 1.44 & 9.52\\
Dejonckheere & Anger & 100 & 7,902 & 10.05 (17.41) & 0--100 & 48.86 & 0.25\\
CES & PHQ-2 & 105 & 4,541 & 1.17 (1.54) & 0--6 & 48.84 & 4.10\\
Marian & Depressed mood & 145 & 7,071 & 1.41 (0.75) & 1--4 & 71.50 & 2.97\\
Marian & Anhedonia & 145 & 7,071 & 1.55 (0.87) & 1--4 & 65.54 & 5.46\\
\bottomrule
\end{tabular}%
}
\begin{flushleft}\footnotesize Summaries describe post-calibration targets, not all scheduled opportunities. Floor and ceiling denote the minimum and maximum admissible score, respectively.\end{flushleft}
\end{table*}

\begin{table*}[!htbp]
\caption{CES demographic information available in the public release.}\label{tab:s12}
\centering\small
\begin{tabular}{@{}llrr@{}}
\toprule
Variable & Level & Count & Percent\\
\midrule
Gender & Female & 71 & 67.6\\
 & Male & 32 & 30.5\\
 & Both & 1 & 1.0\\
 & Missing & 1 & 1.0\\
Race & White & 65 & 61.9\\
 & Asian & 24 & 22.9\\
 & Black & 5 & 4.8\\
 & More than one & 5 & 4.8\\
 & Other/Hispanic & 3 & 2.9\\
 & Other released categories & 2 & 1.9\\
 & Missing & 1 & 1.0\\
\bottomrule
\end{tabular}
\begin{flushleft}\footnotesize Percentages use the 105-participant CES analytic sample. Equivalent harmonised demographic fields were not available in the two openESM analysis files; population descriptions therefore follow their source reports. Categories reproduce the public CES coding and should not be interpreted as investigator-assigned identities.\end{flushleft}
\end{table*}

\begin{table*}[!htbp]
\caption{Window-length sensitivity on common targets after 12 complete reports.}\label{tab:s13}
\centering\small
\begin{tabular}{@{}llrrrr@{}}
\toprule
Dataset & Outcome & $K=4$ & $K=6$ & $K=8$ & $K=12$\\
\midrule
Dejonckheere & Sadness & $-10.10$ & $-11.21$ & $-10.49$ & $-9.23$\\
 & Stress & $-19.38$ & $-19.13$ & $-19.20$ & $-15.71$\\
 & Happiness & $-22.27$ & $-21.10$ & $-19.57$ & $-15.82$\\
 & Relaxation & $-20.08$ & $-19.57$ & $-18.72$ & $-15.09$\\
 & Anger & $-7.80$ & $-6.99$ & $-10.88$ & $-10.54$\\
CES & PHQ-2 & $-29.00$ & $-28.01$ & $-26.52$ & $-22.82$\\
Marian & Depressed mood & $-37.89$ & $-33.87$ & $-29.92$ & $-23.98$\\
 & Anhedonia & $-40.68$ & $-36.40$ & $-31.66$ & $-24.85$\\
\bottomrule
\end{tabular}
\begin{flushleft}\footnotesize Values are participant-balanced relative MAE changes (\%), with negative values favouring the recent $K$-report mean. The common target set contained 7,502 targets per Dejonckheere outcome (100 participants), 4,127 PHQ-2 targets (99 participants), and 6,491 targets per Marian outcome (145 participants). Participant-bootstrap intervals excluded zero except for anger at $K=4$ (95\% CI $-14.62$ to 0.15) and $K=6$ ($-13.52$ to 0.15). Exact intervals for every cell are supplied in the machine-readable source table.\end{flushleft}
\end{table*}

\begin{table*}[!htbp]
\caption{Repeated-origin sensitivity in the four primary outcomes.}\label{tab:s14}
\centering\small
\resizebox{\textwidth}{!}{%
\begin{tabular}{@{}lllrrcc@{}}
\toprule
Dataset & Outcome & Origin & Participants & Targets & MAE change, \% (95\% CI) & Mean age slope (95\% CI)\\
\midrule
Dejonckheere & Sadness & Early & 100 & 7,102 & $-10.95$ ($-16.41$, $-5.39$) & 0.418 (0.094, 0.792)\\
 & & Middle & 100 & 3,828 & $-0.39$ ($-6.96$, 6.46) & 0.457 (0.018, 0.895)\\
 & & Late & 100 & 500 & $-7.67$ ($-20.35$, 6.01) & $-0.059$ ($-0.768$, 0.686)\\
Dejonckheere & Stress & Early & 100 & 7,102 & $-19.95$ ($-24.30$, $-15.67$) & 1.236 (0.616, 1.857)\\
 & & Middle & 100 & 3,828 & $-14.44$ ($-20.22$, $-8.64$) & 0.541 ($-0.083$, 1.202)\\
 & & Late & 100 & 500 & $-11.59$ ($-20.28$, $-2.70$) & $-0.026$ ($-1.128$, 1.084)\\
CES & PHQ-2 & Early & 93 & 3,724 & $-28.90$ ($-36.99$, $-20.59$) & 0.084 (0.034, 0.140)\\
 & & Middle & 93 & 2,114 & $-15.87$ ($-22.03$, $-9.89$) & 0.034 ($-0.007$, 0.077)\\
 & & Late & 93 & 465 & $-13.69$ ($-21.55$, $-5.02$) & 0.048 ($-0.008$, 0.108)\\
Marian & Depressed mood & Early & 145 & 5,911 & $-31.63$ ($-37.88$, $-25.22$) & 0.046 (0.030, 0.064)\\
 & & Middle & 145 & 3,353 & $-26.30$ ($-33.69$, $-19.09$) & 0.019 ($-0.002$, 0.040)\\
 & & Late & 145 & 725 & $-21.42$ ($-33.99$, $-8.11$) & 0.004 ($-0.021$, 0.027)\\
\bottomrule
\end{tabular}%
}
\begin{flushleft}\footnotesize Negative MAE changes favour L8. Positive slopes indicate that the L8 advantage grows with log report-distance from that pseudo-origin. Each origin used eight calibration reports, eight washout reports, and at least five targets. Origins share participants and observations and are not independent replications.\end{flushleft}
\end{table*}

\begin{table*}[!htbp]
\caption{Recall-period and Marian observation-process sensitivities.}\label{tab:s15}
\centering\small
\resizebox{\textwidth}{!}{%
\begin{tabular}{@{}lllrrcc@{}}
\toprule
Dataset & Outcome & Analysis & Participants & Targets & MAE change, \% (95\% CI) & Mean age slope (95\% CI)\\
\midrule
CES & PHQ-2 & Previous-assessment gap $\geq14$ days & 105 & 4,541 & $-23.38$ ($-30.42$, $-16.62$) & 0.105 (0.062, 0.157)\\
Marian & Depressed mood & Unweighted & 145 & 7,071 & $-27.95$ ($-33.18$, $-22.14$) & 0.059 (0.042, 0.079)\\
 & & IPW & 145 & 7,071 & $-28.58$ ($-34.23$, $-22.72$) & 0.060 (0.042, 0.080)\\
Marian & Anhedonia & Unweighted & 145 & 7,071 & $-29.80$ ($-35.15$, $-23.84$) & 0.052 (0.034, 0.073)\\
 & & IPW & 145 & 7,071 & $-30.00$ ($-35.41$, $-24.33$) & 0.054 (0.035, 0.074)\\
\bottomrule
\end{tabular}%
}
\begin{flushleft}\footnotesize Negative MAE changes favour L8. Every CES analytic target met the 14-day minimum relative to the preceding retained report; the 14--30-day rule for visits adjacent in the source data did not prevent longer retained gaps after intermediate exclusions. Marian response probabilities were joined with the true scheduled-opportunity counter and estimated in folds with no participant overlap. Probabilities were clipped to 0.10--0.99; reciprocal inverse-probability weights were then formed and normalised within participant. Bootstrap intervals condition on the fitted response model. The age-gradient values in Tables S14--S15 use the scaling defined for each sensitivity and should not be compared numerically with main-text gradients unless the estimand and scaling coincide. IPW assumes missing at random conditional on observed history and cannot recover unreported states.\end{flushleft}
\end{table*}

\begin{table*}[!htbp]
\caption{Fitted Gaussian AR(1) null challenge for the four primary outcomes.}\label{tab:s16}
\centering\small
\resizebox{\textwidth}{!}{%
\begin{tabular}{@{}lllrrrrr@{}}
\toprule
Dataset & Outcome & Statistic & Observed & Null mean & Null 95\% interval & Standardised distance & One-sided $p$\\
\midrule
Dejonckheere & Sadness & B8$-$L8 MAE & 1.174 & 0.322 & 0.040, 0.649 & 5.49 & 0.0002\\
 & & Age gradient & 0.533 & 0.021 & $-0.121$, 0.170 & 6.98 & 0.0002\\
Dejonckheere & Stress & B8$-$L8 MAE & 3.855 & 0.649 & 0.211, 1.119 & 13.78 & 0.0002\\
 & & Age gradient & 1.568 & 0.042 & $-0.167$, 0.261 & 13.92 & 0.0002\\
CES & PHQ-2 & B8$-$L8 MAE & 0.228 & 0.044 & 0.018, 0.073 & 13.26 & 0.0002\\
 & & Age gradient & 0.110 & 0.006 & $-0.011$, 0.024 & 11.88 & 0.0002\\
Marian & Depressed mood & B8$-$L8 MAE & 0.147 & 0.020 & 0.008, 0.033 & 20.60 & 0.0002\\
 & & Age gradient & 0.059 & 0.002 & $-0.004$, 0.009 & 17.01 & 0.0002\\
\bottomrule
\end{tabular}%
}
\begin{flushleft}\footnotesize Null summaries are based on 5,000 participant-matched stationary Gaussian AR(1) simulations in complete-report index. Observed report counts and age covariates were retained; latent scores were not generated over missed opportunities and transitions did not depend on elapsed time. The CES age-gradient row uses its prespecified minimum-five-target eligibility set (99 participants, 4,523 targets). The minimum corrected one-sided Monte Carlo value is $1/5001=0.0002$. These post hoc tests address only the fitted Gaussian first-order process.\end{flushleft}
\end{table*}

\begin{table*}[!htbp]
\caption{Expanded comparator and participant-SD-standardised performance in primary outcomes.}\label{tab:s17}
\centering\small
\resizebox{\textwidth}{!}{%
\begin{tabular}{@{}llrrrrrrr@{}}
\toprule
Dataset & Outcome & B8 & L8 & EWM8 & Median8 & Mode8 & Online AR(1) & Oracle mean\\
\midrule
Dejonckheere & Sadness & 11.786 & 10.612 & 10.282 & 9.241 & -- & 10.560 & 10.586\\
Dejonckheere & Stress & 21.094 & 17.239 & 16.853 & 16.568 & -- & 17.820 & 17.675\\
CES & PHQ-2 & 0.975 & 0.747 & 0.724 & 0.690 & 0.706 & 0.763 & 0.768\\
Marian & Depressed mood & 0.526 & 0.379 & 0.367 & 0.323 & 0.321 & 0.407 & 0.404\\
\midrule
Dejonckheere & Sadness, standardised & 0.732 & 0.659 & 0.639 & 0.561 & -- & 0.657 & 0.655\\
Dejonckheere & Stress, standardised & 0.957 & 0.768 & 0.752 & 0.734 & -- & 0.798 & 0.785\\
CES & PHQ-2, standardised & 0.932 & 0.731 & 0.713 & 0.645 & 0.660 & 0.752 & 0.750\\
Marian & Depressed mood, standardised & 0.936 & 0.639 & 0.624 & 0.513 & 0.505 & 0.710 & 0.682\\
\bottomrule
\end{tabular}%
}
\begin{flushleft}\footnotesize Upper rows are participant-balanced MAE in outcome units. Standardised rows divide each participant's MAE by that participant's within-person standard deviation before averaging. Lower is better. Oracle mean uses future observations and is a diagnostic rather than a deployable predictor. Full eight-outcome estimates and bootstrap intervals are available in the archived CSV output.\end{flushleft}
\end{table*}

\begin{table*}[!htbp]
\caption{EMA first-day and prompt-timing sensitivities.}\label{tab:s18}
\centering\small
\resizebox{\textwidth}{!}{%
\begin{tabular}{@{}llrrr@{}}
\toprule
Dataset & Outcome & Original & Exclude day 1 and re-anchor & Prompt adjusted\\
\midrule
Dejonckheere & Sadness & $-9.97\%$ & $-10.14\%$ & $-9.96\%$\\
Dejonckheere & Stress & $-18.27\%$ & $-18.10\%$ & $-17.94\%$\\
Marian & Depressed mood & $-27.95\%$ & $-26.47\%$ & $-27.61\%$\\
Marian & Anhedonia & $-29.80\%$ & $-25.63\%$ & $-29.25\%$\\
\bottomrule
\end{tabular}%
}
\begin{flushleft}\footnotesize Values are participant-balanced L8-versus-B8 MAE changes; negative values favour recent history. The original B8 covered a median of two calendar days in Dejonckheere and three in Marian. The exclusion analysis removed all first-day observations before re-establishing the first eight complete reports. Prompt offsets were estimated from all complete reports and are therefore retrospective descriptive adjustments, not strictly past-only forecasts.\end{flushleft}
\end{table*}

\begin{table*}[!htbp]
\caption{Early--late temporal generalisability after participant-SD standardisation.}\label{tab:s19}
\centering\small
\resizebox{\textwidth}{!}{%
\begin{tabular}{@{}llrrr@{}}
\toprule
Dataset & Outcome & Participants & Raw mean-benefit $r$ & Standardised mean-benefit $r$\\
\midrule
Dejonckheere & Sadness & 100 & 0.650 & 0.693\\
Dejonckheere & Stress & 100 & 0.496 & 0.645\\
CES & PHQ-2 & 94 & 0.748 & 0.657\\
Marian & Depressed mood & 134 & 0.683 & 0.793\\
\bottomrule
\end{tabular}%
}
\begin{flushleft}\footnotesize Pearson correlations compare participant ranks in non-overlapping early and late blocks on a common nonzero-variance eligibility set. Standardisation divides update benefit by the participant's within-person standard deviation. Both blocks share B8; the estimates therefore describe repeatability of a shared-reference task statistic, not a formal reliability coefficient or an identified stable psychological trait. Bootstrap intervals excluded zero for all displayed correlations.\end{flushleft}
\end{table*}

\begin{table*}[!htbp]
\caption{Empirical stationary-block model-sensitivity analysis for primary outcomes.}\label{tab:s20}
\centering\small
\resizebox{\textwidth}{!}{%
\begin{tabular}{@{}lllcccc@{}}
\toprule
Dataset & Outcome & Statistic & Observed & Block 8: null mean ($p$) & Block 16: null mean ($p$) & Block 32: null mean ($p$)\\
\midrule
Dejonckheere & Sadness & B8$-$L8 MAE & 1.174 & 0.523 (0.032) & 0.628 (0.064) & 0.760 (0.131)\\
 & & Age gradient & 0.533 & 0.036 (0.001) & 0.041 (0.001) & 0.067 (0.001)\\
Dejonckheere & Stress & B8$-$L8 MAE & 3.855 & 1.117 (0.001) & 1.370 (0.001) & 1.614 (0.001)\\
 & & Age gradient & 1.568 & 0.087 (0.001) & 0.097 (0.001) & 0.109 (0.001)\\
CES & PHQ-2 & B8$-$L8 MAE & 0.228 & 0.064 (0.001) & 0.083 (0.001) & 0.099 (0.001)\\
 & & Age gradient & 0.110 & 0.008 (0.001) & 0.011 (0.001) & 0.014 (0.001)\\
Marian & Depressed mood & B8$-$L8 MAE & 0.147 & 0.029 (0.001) & 0.036 (0.001) & 0.042 (0.001)\\
 & & Age gradient & 0.059 & 0.002 (0.001) & 0.002 (0.001) & 0.003 (0.001)\\
\bottomrule
\end{tabular}%
}
\begin{flushleft}\footnotesize Values come from 1,000 circular geometric stationary-bootstrap draws at each mean block length. One-sided Monte Carlo values use $(1+k)/1001$; 0.001 is the minimum attainable value after correction. The procedure is conditional on each observed empirical sequence and preserves bounded score support, but block lengths were diagnostic choices and are not calibrated psychological timescales. It is a post hoc model-sensitivity analysis, not a universal test of stationarity.\end{flushleft}
\end{table*}

\begin{table*}[!htbp]
\caption{Protocol-locked Corona Health external-extension estimates.}\label{tab:s21}
\centering\footnotesize
{\itshape Panel A: absolute prediction error}\par\smallskip
\begin{tabular}{@{}llrrrr@{}}
\toprule
Outcome & Role & IDs & Targets & B8 MAE & L8 MAE\\
\midrule
PHQ-9 & Primary & 324 & 5,618 & 2.196 & 1.831\\
GAD-7 & Secondary & 327 & 5,663 & 1.946 & 1.619\\
\bottomrule
\end{tabular}
\par\medskip
{\itshape Panel B: paired contrasts}\par\smallskip
\begin{tabular}{@{}lcc@{}}
\toprule
Outcome & Difference (95\% CI) & Change, \% (95\% CI)\\
\midrule
PHQ-9 & 0.365 (0.276, 0.462) & $-16.60$ ($-20.11$, $-13.18$)\\
GAD-7 & 0.327 (0.255, 0.410) & $-16.82$ ($-20.35$, $-13.54$)\\
\bottomrule
\end{tabular}
\par\medskip
{\itshape Panel C: baseline-age gradients and participant-level direction}\par\smallskip
\begin{tabular}{@{}lcc@{}}
\toprule
Outcome & Mean age gradient (95\% CI) & L8 better, IDs\\
\midrule
PHQ-9 & 0.291 (0.168, 0.411)$^{a}$ & 58.0\%\\
GAD-7 & 0.226 (0.141, 0.316)$^{b}$ & 60.9\%\\
\bottomrule
\end{tabular}
\begin{flushleft}\footnotesize Participant-balanced estimates with 2,000 identifier-bootstrap resamples. Positive differences and gradients favour recent history. $^{a}$The PHQ-9 gradient sample contained 232 identifiers and 5,426 targets. $^{b}$The GAD-7 gradient sample contained 232 identifiers and 5,462 targets. The source documentation notes that an anonymised user identifier may change after app reinstallation or a phone change; identifier counts therefore need not equal distinct natural persons. PHQ-9 was frozen as primary; GAD-7 could not replace it based on results.\end{flushleft}
\end{table*}

\begin{table*}[!htbp]
\caption{Corona Health long-gap, ordinal-score, and stationary-process checks.}\label{tab:s22}
\centering\small
\resizebox{\textwidth}{!}{%
\begin{tabular}{@{}lllrrrr@{}}
\toprule
Analysis & Outcome & Method/statistic & Participants & Targets/draws & Estimate & 95\% interval or $p_{\mathrm{MC}}$\\
\midrule
Split at gaps $>30$ days & PHQ-9 & B8$-$L8 MAE & 274 & 4,978 & 0.404 & 0.303, 0.518\\
 & GAD-7 & B8$-$L8 MAE & 276 & 5,019 & 0.351 & 0.270, 0.439\\
Rounded exact accuracy & PHQ-9 & B8 / L8 / Median8 & 324 & 5,618 & 21.0\% / 25.4\% / 29.2\% & Descriptive\\
Rounded within-two accuracy & PHQ-9 & B8 / L8 / Median8 & 324 & 5,618 & 69.6\% / 75.7\% / 76.2\% & Descriptive\\
Rounded exact accuracy & GAD-7 & B8 / L8 / Median8 & 327 & 5,663 & 20.7\% / 25.9\% / 28.1\% & Descriptive\\
Rounded within-two accuracy & GAD-7 & B8 / L8 / Median8 & 327 & 5,663 & 73.5\% / 80.0\% / 79.7\% & Descriptive\\
Gaussian AR(1) null & PHQ-9 & B8$-$L8 MAE & 324 & 5,000 draws & Observed 0.365; null mean 0.175 & Null 0.119, 0.236; $p=0.0002$\\
Gaussian AR(1) null & PHQ-9 & Mean age gradient & 232 & 5,000 draws & Observed 0.291; null mean 0.073 & Null 0.021, 0.129; $p=0.0002$\\
\bottomrule
\end{tabular}%
}
\begin{flushleft}\footnotesize Long-gap sequences were reconstructed after segmentation, so no predictor was calculated across a gap longer than 30 days. Ordinal accuracies were averaged within participant before the group mean. AR(1) intervals are central simulation intervals and Monte Carlo values use $(1+k)/5001$; they are not participant-bootstrap confidence intervals.\end{flushleft}
\end{table*}

\begin{table*}[!htbp]
\caption{Participant-balanced marginal checks for the fitted Gaussian AR(1) reference generators.}\label{tab:s23}
\centering\footnotesize
{\itshape Panel A: location and dispersion}\par\smallskip
\begin{tabular}{@{}llcc@{}}
\toprule
Dataset & Outcome & \shortstack{Mean\\observed / simulated} & \shortstack{Within-person SD\\observed / simulated}\\
\midrule
Dejonckheere & Sadness & 11.63 / 13.87 & 15.13 / 11.87\\
Dejonckheere & Stress & 26.10 / 27.89 & 21.94 / 18.76\\
CES & PHQ-2 & 1.14 / 1.23 & 0.97 / 0.79\\
Marian & Depressed mood & 1.43 / 1.51 & 0.54 / 0.41\\
Corona Health & PHQ-9 & 7.00 / 7.00 & 2.10 / 1.84\\
\bottomrule
\end{tabular}
\par\medskip
{\itshape Panel B: scale boundaries and serial dependence}\par\smallskip
\begin{tabular}{@{}llccc@{}}
\toprule
Dataset & Outcome & \shortstack{Floor fraction\\observed / simulated} & \shortstack{Lag-1\\observed / simulated} & \shortstack{Simulated fraction\\outside scale}\\
\midrule
Dejonckheere & Sadness & 0.443 / 0.256 & 0.179 / 0.150 & 0.000\\
Dejonckheere & Stress & 0.248 / 0.152 & 0.226 / 0.198 & 0.000\\
CES & PHQ-2 & 0.489 / 0.203 & 0.261 / 0.232 & 0.000\\
Marian & Depressed mood & 0.703 / 0.289 & 0.246 / 0.208 & 0.000\\
Corona Health & PHQ-9 & 0.114 / 0.000 & 0.244 / 0.356 & 0.063\\
\bottomrule
\end{tabular}
\begin{flushleft}\footnotesize Values average sequence-level features equally across participants or anonymised user identifiers. Simulated values additionally average 200 diagnostic draws per unit (seed 20260918), separate from the 5,000 draws used for the null contrasts. Development generators clipped continuous Gaussian draws to scale bounds; the Corona PHQ-9 generator did not clip, and 6.3\% of simulated values lay outside 0--27. Neither generator preserved the observed integer score grid. These post hoc checks assess model approximation and are not additional hypothesis tests.\end{flushleft}
\end{table*}

\FloatBarrier
\begingroup
\footnotesize
\setlength{\LTleft}{0pt}
\setlength{\LTright}{0pt}
\begin{longtable}{@{}p{1.8cm}p{3.0cm}p{3.4cm}p{3.4cm}@{}}
\caption{Competing explanations, diagnostic evidence, and remaining uncertainty.}\label{tab:s24}\\
\toprule
Explanation & Diagnostic analysis & What the evidence shows & What remains unresolved\\
\midrule
\endfirsthead
\toprule
Explanation & Diagnostic analysis & What the evidence shows & What remains unresolved\\
\midrule
\endhead
Ordinary temporal persistence & Participant-matched Gaussian AR(1) simulations; empirical stationary-block bootstrap in development datasets & Persistence produces a nonzero recent-history advantage, but fitted-null magnitudes were smaller for the primary outcomes, including external PHQ-9 & Gaussian AR(1) is not a universal stationary model; irregular continuous-time dependence was not identified\\
Unusual study entry & First-day exclusion and later pseudo-origins & Updating often remained useful after re-anchoring, while continuous age gradients weakened & Pseudo-origins were retrospectively selected and reused observations\\
Window or estimator choice & Equal $K=4,6,8,12$ histories; EWM8, median, mode, last value, online AR(1) & The direction was not specific to eight-report arithmetic means; robust summaries helped on bounded scales & No rule was prospectively optimized for clinical decisions\\
Observation process & Opportunity-based inverse-probability weighting in scheduled EMA datasets & Modelled response propensity produced similar estimates & Missing-not-at-random states cannot be recovered from observed history\\
Long irregular gaps & Corona Health segmentation at gaps longer than 30 days & Reconstructing histories within shorter segments preserved the PHQ-9 and GAD-7 directions & A 30-day threshold is a sensitivity choice, not an identified psychological timescale\\
Ordinal and floor effects & Rounded/clipped predictions; recent median or mode & Accuracy and ordinal MAE supported the same direction & Arithmetic score differences do not constitute a full item-response model\\
Stable individual decay trait & Non-overlapping early--late estimates of person-specific gradients & Group-average advantages were clear, but individual gradient ranks did not generalise & Short blocks make individual slopes noisy; absence of evidence is not evidence of no heterogeneity\\
\bottomrule
\end{longtable}
\begin{flushleft}\footnotesize This table separates conclusions supported by the diagnostics from mechanisms that remain unidentified. No row should be read as a causal adjudication.\end{flushleft}
\endgroup

\end{document}
